\begin{lemma}[Partial trace of the identity]
    \label{lem:partial_trace_right_identity}
    \lean{Matrix.partialTraceRight_one}
    \leanok
    For finite-dimensional spaces $H_L$ and $H_R$,
    \[
        \tr_R(\mathbf 1_{L\otimes R})
        =(\dim H_R)\mathbf 1_L.
    \]
\end{lemma}

\begin{proof}
    \leanok
    In product-basis coordinates,
    \[
        \bigl(\tr_R\mathbf 1_{L\otimes R}\bigr)_{ij}
        =\sum_k\delta_{ij}
        =(\dim H_R)\delta_{ij}.
    \]
\end{proof}

\begin{lemma}[Adjoint identity for the right partial trace]
    \label{lem:trace_lift_left_factor_mul}
    \lean{Matrix.trace_leftKroneckerEmbed_mul}
    \leanok
    For every operator $M$ on $H_L$ and every operator $\rho$ on
    $H_L\otimes H_R$,
    \[
        \tr((M\otimes\mathbf 1_R)\rho)
        =\tr(M\tr_R\rho).
    \]
\end{lemma}

\begin{proof}
    \leanok
    Expand both traces in product bases and sum first over the identity entry
    on $H_R$.
\end{proof}

\begin{theorem}[Left absorption by the marginal support projector]
    \label{thm:marginal_support_left_absorption}
    \lean{Matrix.PosSemidef.leftKroneckerEmbed_supportProj_mul_self}
    \leanok
    \uses{lem:trace_lift_left_factor_mul}
    Let $\rho$ be positive semidefinite on $H_L\otimes H_R$, and let $P_L$ be
    the orthogonal projector onto the support of $\tr_R\rho$.
    Then
    \[
        (P_L\otimes\mathbf 1_R)\rho=\rho.
    \]
\end{theorem}

\begin{proof}
    \leanok
    Set $Q_L=\mathbf 1_L-P_L$.  Since
    $Q_L\tr_R\rho=0$, the adjoint identity gives
    \[
        \tr\bigl((Q_L\otimes\mathbf 1_R)\rho\bigr)
        =\tr(Q_L\tr_R\rho)=0.
    \]
    Positivity of $\rho$ then implies
    $(Q_L\otimes\mathbf 1_R)\rho=0$.  Finally,
    $\mathbf 1-(Q_L\otimes\mathbf 1_R)=P_L\otimes\mathbf 1_R$ gives the
    claim.
\end{proof}

\begin{theorem}[Right absorption by the marginal support projector]
    \label{thm:marginal_support_right_absorption}
    \lean{Matrix.PosSemidef.mul_leftKroneckerEmbed_supportProj_self}
    \leanok
    \uses{thm:marginal_support_left_absorption}
    Under the same assumptions,
    \[
        \rho(P_L\otimes\mathbf 1_R)=\rho.
    \]
    For $H_L=H_A\otimes H_X$ and $H_R=H_B$, the two absorption identities are
    precisely the $P_{AX}$ support-projector identities in
    \cite[Appendix~D.2, lines 2228--2235]{Cirac2016MPDO_arXiv}.
\end{theorem}

\begin{proof}
    \leanok
    Apply $(\cdot)^\dagger$ to the left absorption identity.  Since
    $\rho^\dagger=\rho$ and
    $(P_L\otimes\mathbf 1_R)^\dagger=P_L\otimes\mathbf 1_R$,
    \[
        \rho(P_L\otimes\mathbf 1_R)
        =\bigl[(P_L\otimes\mathbf 1_R)\rho\bigr]^\dagger
        =\rho^\dagger
        =\rho.
    \]
\end{proof}

\begin{definition}[Partial trace over the left factor]
    \label{def:partial_trace_left}
    \lean{Matrix.partialTraceLeft, Matrix.partialTraceLeft_apply}
    \leanok
    For an operator $\rho$ on $H_L\otimes H_R$, define
    \[
        (\tr_L\rho)_{r,s}
        =\sum_l \rho_{(l,r),(l,s)}.
    \]
\end{definition}

\begin{theorem}[Hermiticity of the left partial trace]
    \label{thm:partial_trace_left_hermitian}
    \lean{Matrix.partialTraceLeft_isHermitian}
    \leanok
    \uses{def:partial_trace_left}
    If $\rho=\rho^\dagger$, then
    $\tr_L\rho=(\tr_L\rho)^\dagger$.
\end{theorem}

\begin{proof}
    \leanok
    Entrywise Hermiticity gives
    \[
        \overline{(\tr_L\rho)_{r,s}}
        =\sum_k\overline{\rho_{(k,r),(k,s)}}
        =\sum_k\rho_{(k,s),(k,r)}
        =(\tr_L\rho)_{s,r}.
    \]
\end{proof}

\begin{theorem}[Positivity of the left partial trace]
    \label{thm:partial_trace_left_pos}
    \lean{Matrix.PosSemidef.partialTraceLeft}
    \leanok
    \uses{def:partial_trace_left}
    For finite-dimensional $H_L$ and $H_R$, if $\rho\geq 0$, then
    $\tr_L\rho\geq 0$.
\end{theorem}

\begin{proof}
    \leanok
    If $\rho^{(k)}$ is the principal submatrix indexed by
    $\{k\}\times H_R$, then
    \[
        (\tr_L\rho)_{r,s}
        =\sum_k\rho_{(k,r),(k,s)}
        =\sum_k\rho^{(k)}_{r,s}.
    \]
    Each $\rho^{(k)}$ is positive semidefinite, and hence so is their finite
    sum.
\end{proof}

\begin{lemma}[Adjoint identity for the left partial trace]
    \label{lem:trace_lift_right_factor_mul}
    \lean{Matrix.trace_rightKroneckerEmbed_mul}
    \leanok
    \uses{def:partial_trace_left}
    For every operator $M$ on $H_R$,
    \[
        \tr((\mathbf 1_L\otimes M)\rho)
        =\tr(M\tr_L\rho).
    \]
\end{lemma}

\begin{proof}
    \leanok
    Expanding in product bases gives
    \[
        \tr((\mathbf 1_L\otimes M)\rho)
        =\sum_{l,r}\sum_{k,s}\delta_{l,k}M_{rs}
          \rho_{(k,s),(l,r)}
        =\sum_{r,s}M_{rs}\sum_l\rho_{(l,s),(l,r)}
        =\tr(M\tr_L\rho).
    \]
\end{proof}

\begin{theorem}[Left absorption by the right marginal support]
    \label{thm:right_marginal_support_left_absorption}
    \lean{Matrix.PosSemidef.rightKroneckerEmbed_supportProj_mul_self}
    \leanok
    \uses{thm:partial_trace_left_pos,lem:trace_lift_right_factor_mul}
    Let $\rho$ be positive semidefinite.  If $P_R$ is the support projector of
    $\tr_L\rho$, then
    \[
        (\mathbf 1_L\otimes P_R)\rho=\rho.
    \]
\end{theorem}

\begin{proof}
    \leanok
    Set $Q=\mathbf 1_R-P_R$.  The preceding trace identity gives
    \[
        \tr\bigl((\mathbf 1_L\otimes Q)\rho\bigr)
        =\tr\bigl(Q\tr_L\rho\bigr)=0,
    \]
    since $P_R$ is the support projector of $\tr_L\rho$.
    Positivity of $\rho$ then gives
    $(\mathbf 1_L\otimes Q)\rho=0$, and hence
    $(\mathbf 1_L\otimes P_R)\rho=\rho$.
\end{proof}

\begin{theorem}[Right absorption by the right marginal support]
    \label{thm:right_marginal_support_right_absorption}
    \lean{Matrix.PosSemidef.mul_rightKroneckerEmbed_supportProj_self}
    \leanok
    \uses{thm:right_marginal_support_left_absorption}
    Under the same assumptions,
    \[
        \rho(\mathbf 1_L\otimes P_R)=\rho.
    \]
\end{theorem}

\begin{proof}
    \leanok
    Take adjoints in the left absorption identity.
\end{proof}

We now use the explicit tripartite indexing
$H_A\otimes(H_X\otimes H_B)$.  Reassociation identifies it with
$(H_A\otimes H_X)\otimes H_B$ whenever the $AX$ marginal is taken.

\begin{definition}[Tripartite local lifts]
    \label{def:tripartite_local_lifts}
    \lean{TripartiteDecorrelation.liftAX, TripartiteDecorrelation.liftXB,
      TripartiteDecorrelation.liftA, TripartiteDecorrelation.liftB}
    \leanok
    Writing $\widehat M_R$ for the lift of an operator $M_R$ supported on
    $R\in\{AX,XB,A,B\}$, the four lifts to
    $H_A\otimes(H_X\otimes H_B)$ are
    \begin{align*}
        \widehat M_{AX} &= M_{AX}\otimes\mathbf 1_B,
        & \widehat M_{XB} &= \mathbf 1_A\otimes M_{XB},\\
        \widehat M_A &= M_A\otimes\mathbf 1_{XB},
        & \widehat M_B &= \mathbf 1_{AX}\otimes M_B.
    \end{align*}
    The first and fourth formulas use the canonical identification
    $(H_A\otimes H_X)\otimes H_B\cong H_A\otimes(H_X\otimes H_B)$.
\end{definition}

\begin{lemma}[Products and adjoints of tripartite lifts]
    \label{lem:tripartite_local_lift_algebra}
    \lean{TripartiteDecorrelation.liftAX_mul,
      TripartiteDecorrelation.liftXB_mul,
      TripartiteDecorrelation.conjTranspose_liftAX,
      TripartiteDecorrelation.conjTranspose_liftXB,
      TripartiteDecorrelation.conjTranspose_liftA,
      TripartiteDecorrelation.conjTranspose_liftB}
    \leanok
    \uses{def:tripartite_local_lifts}
    The lifts from $AX$ and $XB$ preserve multiplication,
    \[
        \widehat{MN}_{AX}=\widehat M_{AX}\widehat N_{AX},
        \qquad
        \widehat{MN}_{XB}=\widehat M_{XB}\widehat N_{XB},
    \]
    and each of the four lifts from $AX$, $XB$, $A$, and $B$ preserves
    adjoints:
    \[
        (\widehat M_R)^\dagger=\widehat{M^\dagger}_R,
        \qquad R\in\{AX,XB,A,B\}.
    \]
\end{lemma}

\begin{proof}
    \leanok
    For the two-factor lifts,
    \[
        (M_{AX}\otimes\mathbf 1_B)(N_{AX}\otimes\mathbf 1_B)
        =(M_{AX}N_{AX})\otimes\mathbf 1_B,
        \qquad
        (\mathbf 1_A\otimes M_{XB})(\mathbf 1_A\otimes N_{XB})
        =\mathbf 1_A\otimes(M_{XB}N_{XB}).
    \]
    \[
        (M_{AX}\otimes\mathbf 1_B)^\dagger
          =M_{AX}^\dagger\otimes\mathbf 1_B,
        \qquad
        (\mathbf 1_A\otimes M_{XB})^\dagger
          =\mathbf 1_A\otimes M_{XB}^\dagger,
    \]
    and similarly
    \[
        (M_A\otimes\mathbf 1_{XB})^\dagger
          =M_A^\dagger\otimes\mathbf 1_{XB},
        \qquad
        (\mathbf 1_{AX}\otimes M_B)^\dagger
          =\mathbf 1_{AX}\otimes M_B^\dagger,
    \]
    with the tensor factors reassociated to the tripartite space.
\end{proof}

\begin{definition}[Tripartite observable decorrelation]
    \label{def:tripartite_observable_decorrelation}
    \lean{TripartiteDecorrelation.IsObservableDecorrelated}
    \leanok
    \uses{def:tripartite_local_lifts}
    Regions $A$ and $B$ are decorrelated relative to $P$ if, for all Hermitian
    observables $O_A=O_A^\dagger$ and $O_B=O_B^\dagger$,
    \[
        P O_A(\mathbf 1-P)O_B P=0.
    \]
    This is \cite[Definition~D.1, lines 2187--2191]{Cirac2016MPDO_arXiv}.
\end{definition}

\begin{definition}[Complex-linear extension of decorrelation]
    \label{def:tripartite_decorrelation}
    \lean{TripartiteDecorrelation.IsDecorrelated}
    \leanok
    \uses{def:tripartite_local_lifts}
    The complex-linear extension requires
    \[
        P O_A(\mathbf 1-P)O_B P=0
    \]
    for arbitrary complex matrices $O_A$ and $O_B$.
\end{definition}

\begin{theorem}[Observable and complex-linear decorrelation are equivalent]
    \label{thm:tripartite_observable_decorrelation_iff}
    \lean{TripartiteDecorrelation.isObservableDecorrelated_iff}
    \leanok
    \uses{def:tripartite_observable_decorrelation,
      def:tripartite_decorrelation}
    Observable decorrelation holds if and only if its complex-linear extension
    holds.
\end{theorem}

\begin{proof}
    \leanok
    Every complex matrix has the decomposition
    \[
        O=H+\mathrm{i}K,
        \qquad
        H=\frac{O+O^\dagger}{2},
        \qquad
        K=\frac{O-O^\dagger}{2\mathrm{i}},
    \]
    where $H=H^\dagger$ and $K=K^\dagger$.  For analogous decompositions
    $O=H+\mathrm{i}K$ and $O'=H'+\mathrm{i}K'$,
    \begin{align*}
        P O(\mathbf 1-P)O'P
        &=P H(\mathbf 1-P)H'P
          +\mathrm{i}P H(\mathbf 1-P)K'P\\
        &\quad+\mathrm{i}P K(\mathbf 1-P)H'P
          -P K(\mathbf 1-P)K'P.
    \end{align*}
    Each summand vanishes by observable decorrelation applied to its two
    Hermitian components.
\end{proof}

\begin{theorem}[Reverse order of decorrelation]
    \label{thm:tripartite_decorrelation_reverse}
    \lean{TripartiteDecorrelation.IsDecorrelated.reverse}
    \leanok
    \uses{def:tripartite_decorrelation}
    If $P=P^\dagger$ and $A$ and $B$ are decorrelated, then, for arbitrary
    complex matrices $O_A$ and $O_B$,
    \[
        P O_B(\mathbf 1-P)O_A P=0.
    \]
\end{theorem}

\begin{proof}
    \leanok
    Apply decorrelation to $O_A^\dagger$ and $O_B^\dagger$, and take adjoints.
\end{proof}

\begin{definition}[Tripartite reduced operators]
    \label{def:tripartite_marginals}
    \lean{TripartiteDecorrelation.marginalAX,
      TripartiteDecorrelation.marginalXB}
    \leanok
    The two reduced operators are
    \[
        \rho_{AX}=\tr_B P,
        \qquad
        \rho_{XB}=\tr_A P.
    \]
\end{definition}

\begin{theorem}[Matrix-unit expansions of the reduced operators]
    \label{thm:tripartite_marginal_matrix_units}
    \lean{TripartiteDecorrelation.lift_marginalAX_eq_sum,
      TripartiteDecorrelation.lift_marginalXB_eq_sum}
    \leanok
    \uses{def:tripartite_local_lifts,def:tripartite_marginals}
    For matrix units on $A$ and $B$,
    \[
        \mathbf 1_A\otimes\rho_{XB}
          =\sum_{a,k}E^A_{ak}P E^A_{ka},
        \qquad
        \rho_{AX}\otimes\mathbf 1_B
          =\sum_{b,k}E^B_{bk}P E^B_{kb}.
    \]
\end{theorem}

\begin{proof}
    \leanok
    For the $XB$ marginal, evaluation in the product basis gives
    \[
        \bigl(E^A_{ak}P E^A_{ka}\bigr)_{(a_0,r),(a_1,s)}
        =\delta_{a,a_0}\delta_{a,a_1}P_{(k,r),(k,s)}.
    \]
    Hence
    \[
        \sum_{a,k}\bigl(E^A_{ak}P E^A_{ka}\bigr)_{(a_0,r),(a_1,s)}
        =\delta_{a_0,a_1}\sum_kP_{(k,r),(k,s)}
        =\bigl(\mathbf 1_A\otimes\rho_{XB}\bigr)_{(a_0,r),(a_1,s)}.
    \]
    The $AX$ identity follows symmetrically.
\end{proof}

\begin{theorem}[Decorrelation for the reduced operators]
    \label{thm:tripartite_marginal_complementary_product_zero}
    \lean{TripartiteDecorrelation.IsDecorrelated.marginals_mul_complement_mul_eq_zero}
    \leanok
    \uses{def:tripartite_decorrelation,
      thm:tripartite_marginal_matrix_units}
    If $A,B$ are decorrelated, then
    \[
        \widehat\rho_{XB}(\mathbf 1-P)\widehat\rho_{AX}=0.
    \]
\end{theorem}

\begin{proof}
    \leanok
    Substitute the matrix-unit expansions.  Every summand has the form
    \[
        E^A_{ak}PE^A_{ka}(\mathbf 1-P)E^B_{bl}PE^B_{lb}
        =E^A_{ak}\bigl[PE^A_{ka}(\mathbf 1-P)E^B_{bl}P\bigr]E^B_{lb}
        =0,
    \]
    where the bracketed factor vanishes by decorrelation.
\end{proof}

\begin{theorem}[Reverse reduced-operator decorrelation]
    \label{thm:tripartite_reverse_marginal_complementary_product_zero}
    \lean{TripartiteDecorrelation.IsDecorrelated.reverse_marginals_mul_complement_mul_eq_zero}
    \leanok
    \uses{thm:tripartite_decorrelation_reverse,
      thm:tripartite_marginal_matrix_units}
    If $P=P^\dagger$ and $A,B$ are decorrelated, then
    \[
        \widehat\rho_{AX}(\mathbf 1-P)\widehat\rho_{XB}=0.
    \]
\end{theorem}

\begin{proof}
    \leanok
    Substitute the matrix-unit expansions.  Every summand has the form
    \[
        E^B_{bl}PE^B_{lb}(\mathbf 1-P)E^A_{ak}PE^A_{ka}
        =E^B_{bl}\bigl[PE^B_{lb}(\mathbf 1-P)E^A_{ak}P\bigr]E^A_{ka}
        =0,
    \]
    where the bracketed factor vanishes by reverse-order decorrelation.
\end{proof}

\begin{definition}[Tripartite marginal support projectors]
    \label{def:tripartite_marginal_supports}
    \lean{TripartiteDecorrelation.supportAX, TripartiteDecorrelation.supportXB}
    \leanok
    \uses{def:partial_trace_left}
    For a Hermitian operator $P$, let $P_{AX}$ and $P_{XB}$ be the support projectors of
    $\tr_B P$ and $\tr_A P$, respectively, as in
    \cite[Appendix~D.2, lines 2225--2229]{Cirac2016MPDO_arXiv}.  Write
    $\widehat P_{AX}$ and $\widehat P_{XB}$ for their lifts to the tripartite
    space.
\end{definition}

\begin{theorem}[Decorrelation for the two marginal supports]
    \label{thm:tripartite_complementary_product_zero}
    \lean{TripartiteDecorrelation.IsDecorrelated.supports_mul_complement_mul_eq_zero,
      TripartiteDecorrelation.IsDecorrelated.reverse_supports_mul_complement_mul_eq_zero}
    \leanok
    \uses{def:tripartite_marginal_supports,
      thm:tripartite_marginal_complementary_product_zero,
      thm:tripartite_reverse_marginal_complementary_product_zero}
    If $P=P^\dagger$ and $A,B$ are decorrelated, then
    \[
        \widehat P_{XB}(\mathbf 1-P)\widehat P_{AX}=0,
        \qquad
        \widehat P_{AX}(\mathbf 1-P)\widehat P_{XB}=0.
    \]
\end{theorem}

\begin{proof}
    \leanok
    For $g(0)=0$ and $g(x)=x^{-1}$ when $x\ne0$, functional calculus gives
    \[
        P_{AX}=g(\rho_{AX})\rho_{AX}=\rho_{AX}g(\rho_{AX}),
        \qquad
        P_{XB}=g(\rho_{XB})\rho_{XB}=\rho_{XB}g(\rho_{XB}).
    \]
    Hence
    \[
        \widehat P_{XB}(\mathbf 1-P)\widehat P_{AX}
        =\widehat{g(\rho_{XB})}
          \bigl[\widehat\rho_{XB}(\mathbf 1-P)\widehat\rho_{AX}\bigr]
          \widehat{g(\rho_{AX})}
        =0,
    \]
    and, in the reverse order,
    \[
        \widehat P_{AX}(\mathbf 1-P)\widehat P_{XB}
        =\widehat{g(\rho_{AX})}
          \bigl[\widehat\rho_{AX}(\mathbf 1-P)\widehat\rho_{XB}\bigr]
          \widehat{g(\rho_{XB})}
        =0.
    \]
    % Local fix for source lines 2236--2252: see
    % docs/paper-gaps/cpsv16_decorrelation_slice_expansion_fix.tex.
\end{proof}

\begin{theorem}[Absorption by the lifted $XB$ support projector]
    \label{thm:tripartite_xb_support_absorption}
    \lean{TripartiteDecorrelation.liftSupportXB_mul_self,
      TripartiteDecorrelation.mul_liftSupportXB_self}
    \leanok
    \uses{def:tripartite_local_lifts,def:tripartite_marginal_supports,
      thm:right_marginal_support_left_absorption,
      thm:right_marginal_support_right_absorption}
    For $P\geq0$,
    \[
        \widehat P_{XB}P=P=P\widehat P_{XB}.
    \]
\end{theorem}

\begin{proof}
    \leanok
    For the bipartition $A|XB$, the two support-absorption identities give
    \[
        (\mathbf 1_A\otimes P_{XB})P
        =P
        =P(\mathbf 1_A\otimes P_{XB}).
    \]
\end{proof}

\begin{theorem}[Absorption by the lifted $AX$ support projector]
    \label{thm:tripartite_ax_support_absorption}
    \lean{TripartiteDecorrelation.liftSupportAX_mul_self,
      TripartiteDecorrelation.mul_liftSupportAX_self}
    \leanok
    \uses{def:tripartite_local_lifts,def:tripartite_marginal_supports,
      thm:marginal_support_left_absorption,
      thm:marginal_support_right_absorption}
    For $P\geq0$,
    \[
        \widehat P_{AX}P=P=P\widehat P_{AX}.
    \]
\end{theorem}

\begin{proof}
    \leanok
    After reassociating the tripartite space as $AX|B$, the two
    support-absorption identities give
    \[
        (P_{AX}\otimes\mathbf 1_B)P
        =P
        =P(P_{AX}\otimes\mathbf 1_B).
    \]
\end{proof}

\begin{theorem}[Product of the two marginal support projectors]
    \label{thm:tripartite_support_products}
    \lean{TripartiteDecorrelation.supportProducts_eq}
    \leanok
    \uses{def:tripartite_marginal_supports,
      thm:tripartite_xb_support_absorption,
      thm:tripartite_ax_support_absorption,
      thm:tripartite_complementary_product_zero}
    If $P=P^\dagger=P^2$ and $A$ and $B$ are decorrelated relative to $P$, then
    \[
        \widehat P_{XB}\widehat P_{AX}
        =P
        =\widehat P_{AX}\widehat P_{XB}.
    \]
\end{theorem}

\begin{proof}
    \leanok
    Inserting $\mathbf 1=P+(\mathbf 1-P)$ between the two support projectors
    gives
    \[
        \widehat P_{XB}\widehat P_{AX}
        =\widehat P_{XB}P\widehat P_{AX}
          +\widehat P_{XB}(\mathbf 1-P)\widehat P_{AX}
        =P+0
        =P.
    \]
    Here the second equality uses the support-absorption identities and
    Theorem~\ref{thm:tripartite_complementary_product_zero}.  The same
    calculation with the two factors reversed gives
    $\widehat P_{AX}\widehat P_{XB}=P$.
\end{proof}
